\reading{Bikram's Yoga College v. Evolation Yoga}
\readingcite{803 F.3d 1032 (9th Cir. 2015)}


\opinion \textsc{Wardlaw}, Circuit Judge.

We must decide whether a sequence of twenty-six yoga poses and two
breathing exercises developed by Bikram Choudhury and described in his
1979 book, \textit{Bikram's Beginning Yoga Class}, is entitled to
copyright protection.\ldots

\readinghead{I. Factual and Procedural History}

The Indian practice and philosophy of yoga date back thousands of years.
Derived from ancient Hindu scriptures, including the Bhagavad
Gita, the practice of yoga teaches students to attain spiritual
fulfillment  through control of the mind and body.
Yoga has evolved into a diverse set of spiritual,
philosophical, and physical disciplines. Some students practice yoga to
transcend the physical body and unite with divine powers; others focus
on improving strength, flexibility, and overall physical fitness.\ldots

%The history of yoga in the United States reflects its wide-ranging
%appeal. Some of yoga's first American adherents included
%nineteenth-century transcendentalists, such as Henry David Thoreau and
%Ralph Waldo Emerson, who were fascinated by yoga's approach to achieving
%enlightenment. In the early twentieth century, yoga grew more popular
%as scientists and physicians began to study the physical benefits of the
%practice. These physical benefits caught the attention of Hollywood
%celebrities, including Gloria Swanson, Greta Garbo, and Marilyn Monroe,
%who embraced yoga as a tool to fight illness and aging.
%By the 1960s, Americans increasingly turned to yoga as a
%``non-religious, decidedly unspiritual'' form of physical exercise.

In 1971, Bikram Choudhury, the ``self-proclaimed `Yogi to the stars,'\,''
arrived in Beverly Hills, California. He soon became a central
figure in the growing popularity of yoga in the United States. Born and
raised in Calcutta, India, Choudhury began studying yoga at age four and
learned hundreds of traditional Hatha yoga ``asanas,'' or individual
poses. Hatha yoga places particular emphasis on the physical components
of yoga. Choudhury developed a sequence of twenty-six asanas and two
breathing exercises, arranged in a particular order, which he calls the
``Sequence.''\ldots
% Choudhury opened his own studio, where he began offering
%``Bikram Yoga'' classes. In a Bikram Yoga class, the Sequence is practiced
%over the course of ninety minutes, to a series of instructions (the
%``Dialogue''), in a room heated to 105 degrees Fahrenheit to simulate
%Choudhury's native Indian climate.
%
%Choudhury popularized the Sequence by marketing the many health and
%fitness benefits it provides. Choudhury informs prospective students
%that his ``system of Hatha Yoga is capable of helping you avoid, correct,
%cure, heal, or at least alleviate the symptoms of almost any illness or
%injury.'' He claims that he developed the Sequence after ``many of years
%of research and verification\ldots using modern medical measurement
%techniques.'' He tells reporters that he extended the careers of
%professional athletes, including Kareem Abdul-Jabbar and John McEnroe.
%This message has resonated with an American audience: as the complaint
%in this action explains, ``[p]ublic demand for Bikram Yoga classes grew
%steadily once Bikram Yoga participants realized that Bikram's unique
%yoga style and method offered them tremendous physical, mental and other
%benefits.''

In 1979, Choudhury published the book \textit{Bikram's Beginning Yoga Class},
which includes descriptions, photographs, and drawings of the
Sequence's twenty-six poses and two breathing exercises. Choudhury
registered the book with the U.S. Copyright Office in 1979. In 2002, he
also registered the ``compilation of exercises'' contained in the book,
using a supplementary  registration form that referenced back to the 1979 book.

%In 1994, Choudhury introduced the ``Bikram Yoga Teacher Training
%Course.'' In 2002 and 2005, respectively, Mark Drost and Zefea Samson
%enrolled in and successfully completed the three-month Bikram Yoga
%Teacher Training course. In 2009, Drost and Samson founded Evolation
%Yoga, LLC. Evolation Yoga offers several types and styles of yoga,
%including ``hot yoga,'' which is similar to ``Bikram's Basic Yoga System.''
%Evolation acknowledges that hot yoga ``includes 26 postures and two
%breathing exercises and is done for 90 minutes, accompanied by a series
%of oral instructions, in a room heated to approximately 105 degrees
%Fahrenheit.''
%
%On July 1, 2011, Choudhury and Bikram's Yoga College of India, L.P.
%(``Choudhury'') filed a complaint in the Central District of California
%alleging, \textit{inter alia}, that defendants Evolation Yoga, LLC, Mark Drost,
%and Zefea Samson
%(``Evolation'') infringed ``Bikram's Copyrighted Works through substantial
%use of Bikram's Copyrighted Works in and as part of Defendants' offering
%of yoga classes.'' On November 12, 2012, Evolation moved for partial
%summary judgment as to Choudhury's claim of copyright infringement of
%the ``Sequence.'' The district court granted Evolation's motion, ruling
%that the ``Sequence is a collection of facts and ideas'' that is not
%entitled to copyright protection. The parties settled all remaining
%claims against each other, and Choudhury timely appealed as to the
%``Sequence.''

[Two of Choudhury's students founded Evolation
Yoga, which offered similar yoga programs. Choudhury sued Evolation for
copyright infringement.]

%\readinghead{II. Standard of Review}
%
%``We review de novo a district court's grant of partial summary judgment, and
%may affirm on any ground supported by the record.'' ``After `viewing the
%evidence in the light most favorable to the nonmoving party,' we determine
%`whether there are any genuine issues of material fact and whether the district
%court correctly applied the relevant substantive law.'\,''

\readinghead{III. Discussion}

Though Choudhury emphasizes the aesthetic attributes of the
Sequence's ``graceful flow,'' at bottom, the Sequence is an idea, process,
or system designed to improve health. Copyright protects only the
expression of this idea---the words and pictures used to describe the
Sequence---and not the idea of the Sequence itself.\ldots
% Because the Sequence
%is an unprotectable idea, it is also ineligible for copyright
%protection as a ``compilation'' or ``choreographic work.'' The district
%court properly granted partial summary judgment in favor of Evolation
%because the Sequence is not a proper subject of copyright.



\readinghead{\textit{A. The Sequence Is an Unprotectable Idea.}}



[17 U.S.C. \S] 102(a) of the Copyright Act of 1976 sets forth the proper
subjects of copyright protection.
Section 102(b) expressly excludes protection for ``any idea, procedure,
process, system, method of operation, concept, principle, or discovery,
regardless of the form in which it is described, explained, illustrated,
or embodied in such work.'' Section 102(b) codifies
the ``idea/expression dichotomy,'' under which ``every idea, theory, and
fact in a copyrighted work becomes instantly available for public
exploitation at the moment of publication.'' \textit{Golan v. Holder}, 565 U.S.
302, 328 (2012)\ldots.

The idea/expression dichotomy has two constitutional foundations: the
Copyright Clause and the First Amendment. Under the Copyright Clause,
``[t]he primary objective of copyright is not to reward the labor of
authors, but `[t]o promote the Progress of Science and useful Arts.'\,''
Thus, ``[t]he `constitutional
command'\ldots is that Congress, to the extent it enacts copyright laws at
all, create a `system' that `promote[s] the Progress of Science.'\,''
``To this end, copyright assures authors the right to their original
expression, but encourages others to build freely upon the ideas and
information conveyed by a work.''
At the same time, the idea/expression dichotomy ``strike[s] a
definitional balance between the First Amendment and the Copyright Act
by permitting free communication of facts while still protecting an
author's expression.''\edfootnote{The omitted citations in this paragraph are
primarily to Supreme Court copyright cases.}

In \textit{Baker v. Selden}, 101 U.S. 99 (1879),
the Supreme Court addressed the protection copyright law provided to a
book, a classic subject of copyright protection, explaining a \textit{system} of
book-keeping. The Court held that the book's expression of the book-keeping
system was protected, but the system of book-keeping itself was not entitled to
copyright protection. The Court explained:
\begin{quote}
The description of the art\edfootnote{In the late 1800s, ``art'' meant a
craft or trade. This usage remains in phrases like ``the art of woodworking'' or
``the dark arts.''} in a book, though entitled to the
benefit of copyright, lays no foundation for an exclusive claim to the
art itself. The object of the one is explanation; the object of the
other is use. The former may be secured by copyright. The latter can
only be secured, if it can be secured at all, by letters-patent.
\end{quote}

Following \textit{Baker}, and recognizing this vital distinction between ideas
and expression, courts have routinely held that the  copyright for a work
describing how to perform a process does not extend to the process itself.
[As support, the court cited cases dealing with meditation exercises, recipes in
a cookbook, and rules for roller-skating races.] 
%In
%\textit{Palmer v. Braun}, 287 F.3d 1325 (11th Cir. 2002),
%for example, the Eleventh Circuit held that meditation exercises
%described in a copyrighted manual on exploring the consciousness were ``a
%process'' unentitled to copyright protection. The
%court explained that the ``exercises, while undoubtedly the product of
%much time and effort, are, at bottom, simply a process for achieving
%increased consciousness. Such processes, even if original, cannot be
%protected by copyright.'' Similarly, in \textit{Publications International, Ltd.
%v. Meredith Corp.}, 88 F.3d 473 (7th Cir. 1996),
%the Seventh Circuit held that recipes contained in a copyrighted
%cookbook are not entitled to copyright protection, for they merely
%``describe a procedure by which the reader may produce many dishes,'' and
%``there can be no monopoly in the copyright sense in the ideas for
%producing certain foodstuffs.'' Finally, in \textit{Seltzer v. Sunbrock}, 22
%F. Supp. 621 (S.D. Cal. 1938), which predates the Copyright Act of 1976 but
%applies
%\textit{Baker},
%the court held that the copyright in a manual describing how to
%organize roller-skating races does not extend to the rules for the races
%themselves. The court explained, ``[w]hat [the
%author] really composed was a description of a system for conducting
%races on roller skates. A system, as such, can never be copyrighted. If
%it finds any protection, it must come from the patent laws.''

Here, we must similarly determine not the validity of a copyright but rather its
scope.
Does Choudhury's copyright protection for his 1979 book extend to the
Sequence itself? Under the fundamental tenets of copyright law and
consistent with the precedents discussed above, the answer is no.

As Choudhury describes it, the Sequence is a ``system'' or a ``method''
designed to ``systematically work every part of the body, to give all
internal organs, all the veins, all the ligaments, and all the muscles
everything they need to maintain optimum health and maximum function.''\ldots
%In \textit{Bikram's Beginning Yoga Class}, Choudhury explains that he
%``arrived at the sequence of postures'' after ``[researching] the diseases
%and the postures and after many years of research and verification\ldots
%using modern medical measurement techniques.'' The book tells readers
%that ``Bikram's twenty-six exercises systematically move fresh,
%oxygenated blood to one hundred percent of your body, to each organ and
%fiber, restoring all systems to healthy working order, just as Nature
%intended.''
This text promises readers that Choudhury's ``system of
Hatha Yoga is capable of helping you avoid, correct, cure, heal, or at
least alleviate the symptoms of almost any illness or injury.''

[The court reviewed many statements from Choudhury's book explaining the health
benefits of the Sequence, including statements about the
health importance of following the sequence in ``strict order.'']

%Also illuminating is Choudhury's spoken Dialogue, which accompanies
%the Sequence. Before the Sequence's first breathing exercise, for
%example, the instructor tells students, ``[The exercise] is good for the
%lungs and respiratory system. This exercise expands your lungs to their
%maximum expansion capacity. And it improves the elasticity of your
%lungs.'' Before the twelfth pose, the instructor explains:
%\begin{quote}
%Every exercise in the world you do, you burn energy/calories
%like driving a car burns gas. The tank is empty, you need to fill it up
%again. Hatha Yoga class is a gas station, it is the only place in the
%world where you gain energy instead of burning energy. Asana is the only
%\textit{natural} physical activity in the world because it is scientific [and]
%with the help of science, we can explain nature.
%\end{quote}
%
%An essential element of this ``system'' is the order in which the yoga poses and breathing exercises are arranged. \textit{Bikram's Beginning Yoga Class}
%instructs readers, ``Do the poses in the strict order given in this
%book. Nothing about Bikram's Beginning Yoga Class is haphazard. It is
%designed to scientifically warm and stretch muscles, ligaments, and
%tendons in the order in which they should be stretched.''
%For instance, Choudhury explains, ``Camel Pose (Ustrasana)
%stretches the abdomen and compresses the spine; so for the next posture,
%I chose the Rabbit Pose (Sasangasana), which does the converse:
%stretches the back and compresses the abdomen.'' One \textit{Yoga Journal}
%article explains that ``[a]ccording to Bikram, each posture in his
%series forms the perfect basis for the next, warming and stretching the
%appropriate muscles, ligaments and tendons.''

Choudhury thus attempts to secure copyright protection for a healing
art: a system designed to yield physical benefits and a sense of
well-being. Simply put, this attempt is precluded by copyright's
idea/expression dichotomy, codified by Section 102(b).
As the Supreme
Court explained in \textit{Baker}, ``Certain mixtures are found to be of
great value in the healing art. If the discoverer writes and publishes a
book on the subject (as regular physicians generally do), he gains no
exclusive right to the manufacture and sale of the medicine; he gives
that to the public.''\ldots
%Thus, for example, the copyright for a book describing how to perform a
%complicated surgery does not give the holder the exclusive right to
%perform the surgery. Like the series of movements a surgeon makes, the
%Sequence is, as Choudhury tells readers, a method designed to ``cure,
%heal, or at least alleviate'' physical injuries and illness. Monopoly
%protection for such a method ``can only be secured, if it can be secured
%at all, by letters-patent.''
In light of \textit{Baker} and its progeny, Choudhury's healing methodology
is not eligible for protection by copyright. Indeed, if it is entitled
to protection at all, that protection is more properly sought through
the patent process.\ldots

%That the Sequence may produce spiritual and psychological benefits
%makes it no less an idea, system, or process and no more amenable to
%copyright protection. Choudhury's personal declaration explains that the
%Sequence offers ``spiritual benefits'' to his students and ``lead[s] to a
%general sense of peace and well-being that is undoubtedly of benefit to
%all of us.'' Like the meditation exercises designed to achieve greater
%consciousness in \textit{Braun}, 287 F.3d at 1334,
%the Sequence sets forth a method to attain identifiable, if spiritual
%and psychological, results: a ``sense of well-being'' and ``boundless
%energy.'' As such, it falls within the Copyright Act's definition of an idea,
%process, or system excluded from copyright protection.

Choudhury contends that the Sequence's arrangement of postures is
``particularly beautiful and graceful.'' But beauty is not a basis for
copyright protection. The performance of many ideas, systems, or
processes may be beautiful: a surgeon's intricate movements, a
book-keeper's careful notations, or a baker's kneading might each
possess a certain grace for at least some viewers. Indeed, from
Vermeer's milkmaid to Lewis Hine's power house mechanic, the individual
engrossed in a process has long attracted artistic attention. But the
beauty of the process does not permit one who describes it to gain,
through copyright, the monopolistic power to exclude all others from
practicing it. This is true even where, as here, the process was
conceived with at least some aesthetic considerations in mind. Just as
some steps in a recipe may reflect no more than the author's belief that
a particular ingredient is beautiful or that a particular cooking
technique is impressive to watch and empowering to practice, some
elements in Choudhury's Sequence may reflect his aesthetic preferences.
Yet just like the recipe, the Sequence remains unprotectable as a
process the design of which primarily reflects function, not expression.

In drawing the ``difficult'' line between idea and expression in this
case, we are mindful of the ``guiding consideration'' of the
idea/expression dichotomy: ``the preservation of the balance between
competition and protection reflected in the patent and copyright laws.''  As in
\textit{Baker}, the ``object'' of the book \textit{Bikram's Beginning Yoga
Class} is ``explanation'': it tells readers how to perform the Sequence and
encourages them to try it. The introduction to
\textit{Bikram's Beginning Yoga Class}, for example, urges the audience to: (i)
``turn to the Contents page,'' (ii) ``read  through the book,'' (iii) ``build
gradually,'' and (iv) ``do the poses in the strict order given in this book.''
Like a book explaining
``Book-keeping Simplified,'' \textit{Bikram's Beginning Yoga
Class} sets out to ``communicate to the world the useful knowledge which it
contains.''
It invites readers to practice the method it describes. ``But
this object would be frustrated if the knowledge could not be used
without incurring the guilt of piracy of the book.'' Consumers
would have little reason to buy Choudhury's book if Choudhury held a
monopoly on the practice of the very activity he sought to popularize.
Rather than ``stimulat[ing] artistic creativity for the general public
good,'' copyright protection for the Sequence would prevent the public
from engaging with Choudhury's idea and building upon it. 

%\readinghead{\textit{B. The Sequence Is Not a Copyrightable Compilation.}}
%
%Choudhury contends that the Sequence is entitled to copyright
%protection as a ``compilation.'' Specifically, Choudhury claims that the
%Sequence qualifies for copyright protection because his ``selection,
%coordination, and arrangement'' of twenty-six poses and two breathing
%exercises create a coherent and expressive composition. The district
%court correctly rejected this argument.
%
%The Copyright Act identifies compilations as a proper subject of
%copyright. Section 103 of the Copyright Act provides that ``[t]he subject
%matter of copyright as specified in section 102 includes compilations.''
%A ``compilation'' is ``a work formed by the collection
%and assembling of preexisting materials or of data that are selected,
%coordinated, or arranged in such a way that the resulting work as a
%whole constitutes an original work of authorship.'' It
%essential to recognize, however, that Section 103 complements Section
%102. Thus, while a compilation may be eligible for copyright protection,
%it must nevertheless satisfy the requirements of Section 102. A
%compilation must, in other words, represent an ``original work[] of
%authorship,'' and ``[i]n no case'' may copyright protection ``extend to any
%idea, procedure, process, [or] system.'' The
%availability of copyright protection for compilations, therefore, does
%not eliminate Section 102's categorical bar on copyright protection for
%ideas.
%
%The Supreme Court addressed the relationship between these ``two
%well-established propositions''---that compilations are eligible for
%copyright but facts and ideas are not---in \textit{Feist}, 499 U.S. 340 (1991).
%In \textit{Feist},
%the Court considered whether the collection of names, towns, and
%telephone numbers in a telephone directory is eligible for copyright
%protection as a compilation. The Court held that ``[a] factual
%compilation is eligible for copyright if it features an original
%selection or arrangement of facts, but the copyright is limited to the
%particular selection or arrangement. In no event may copyright extend to
%the facts themselves.''
%
%By claiming copyright protection for the Sequence as a compilation,
%Choudhury misconstrues the scope of copyright protection for
%compilations. As we have explained, the Sequence is an idea, process, or
%system; therefore, it is not eligible for copyright protection. That
%the Sequence may possess many constituent parts does not transform it
%into a proper subject of copyright protection. Virtually any process or
%system could be dissected in a similar fashion. \textit{Baker}'s examples of
%``how-to'' treatises are instructive: ``A treatise
%on\ldots the construction and use of ploughs, or watches, or churns[,]\ldots or
%on the mode of drawing lines to produce the effect of
%perspective'' would likely list the steps necessary to perform the
%process it describes.
%The watchmaking treatise's author could not claim a copyright in the
%process of making a watch, however, by breaking down the process into
%multiple steps and labeling it a ``compilation.'' Recipes further
%illustrate the point: a cake recipe could be viewed as a ``compilation''
%of carefully arranged and selected steps---which may, of course, reflect
%the personal preferences and tastes of the recipe's author---yet the
%recipe would remain, in most instances, a process that is not eligible
%for copyright protection. \textit{See} \textit{Meredith}, 88 F.3d at 480-81.
%Likewise, Choudhury cannot obtain copyright protection for the Sequence
%as a compilation by separately identifying the poses and breathing
%exercises it contains.
%
%Moreover, according to Choudhury himself, the medical and functional
%considerations at the heart of the Sequence compel the very selection
%and arrangement of poses and breathing exercises for which he claims
%copyright protection. According to \textit{Bikram's Beginning Yoga Class},
%the ``strict order'' of the poses ``is designed to scientifically warm and
%stretch muscles, ligaments, and tendons in the order in which they
%should be stretched.'' \textit{Bikram's Beginning Yoga Class, supra}, at
%xi. Read in the light most favorable to Choudhury, the record
%demonstrates that the overarching reason for the organization of the
%poses and breathing exercises in the Sequence is to further the basic
%goals of the method: to attain ``[p]roper weight, muscle tone, glowing
%complexion, boundless energy, vibrant good health, and a sense of
%well-being.'' The Sequence's composition renders it more
%effective as a process or system, but not any more suitable for
%copyright protection as an original work of authorship.
%
%It makes no difference that similar results could be achieved through
%a different organization of yoga poses and breathing exercises.
%Choudhury argues that he could have chosen from ``hundreds of postures''
%and ``countless arrangements of these postures'' in developing the
%Sequence. But the possibility of attaining a particular end through
%multiple different methods does not render the uncopyrightable a proper
%subject of copyright. Though it may be one of many possible yoga
%sequences capable of attaining similar results, the Sequence is
%nevertheless a process and is therefore ineligible for copyright
%protection.
%
%\readinghead{\textit{C. The Sequence Is Not a Copyrightable Choreographic
%Work.}}
%
%
%
%The district court properly rejected Choudhury's argument that the
%Sequence is entitled to copyright protection as a choreographic work.
%The 1976 Copyright Act extended protection to ``pantomimes and
%choreographic works,'' which were previously not copyrightable. In
%1986, the Second Circuit observed that ``[e]xplicit federal copyright
%protection for choreography is a fairly recent development, and the
%scope of that protection is an uncharted area of the law.'' This remains true
%today.
%
%The parties debate the meaning of the term ``choreography,'' which we
%have not yet defined in the copyright context. Nor did Congress define
%the term ``choreographic work[],'' apparently because its meaning was
%``fairly settled.'' The legislative
%history does explain, however, that it is not ``necessary to specify that
%`choreographic works' do not include social dance steps and simple
%routines.'' The Second Circuit has relied on the
%Compendium of Copyright Office Practices as persuasive authority and
%concluded that ``[c]horeography represents a related series of dance
%movements and patterns organized into a coherent whole.'' The \textit{Compendium
%II} defines ``dance'' as ``static and kinetic successions of bodily movement in
%certain rhythmic and spatial relationships.''
%The ``dance movements,'' according to the \textit{Compendium II}, ``must be more
%than mere exercises, such as `jumping jacks' or walking steps.''
%Finally, the \textit{Compendium II}
%explains that choreography is ``usually intended to be accompanied by
%music'' but ``need not tell a story'' and need not be presented ``before an
%audience.''
%
%In this case, we need not decide whether to adopt the Copyright
%Office's definition of ``choreographic work'' or fashion another on our
%own because all categories of works eligible for copyright protection,
%including choreographic works, are subject to the critical requirements
%and limitations of Section 102. The beauty of this section is that it
%allows for the possibility that the term ``original work of authorship''
%may, as
%it has, evolve and encompass new forms of expression that, like
%choreography, are not easily reduced to neat definitions. ``Congress
%recurrently adjusts copyright law to protect categories of works once
%outside the law's compass.''
%Yet as Congress has responded to new technologies and evolving
%understandings of creative expression, the idea/expression dichotomy has
%remained firmly in place. This dichotomy, as this case illustrates,
%polices the uncertain boundaries of copyrightable subject matter.
%
%The Sequence is not copyrightable as a choreographic work for the
%same reason that it is not copyrightable as a compilation: it is an
%idea, process, or system to which copyright protection may ``[i]n no
%case'' extend. We recognize that the Sequence may
%involve ``static and kinetic successions of bodily movement in certain
%rhythmic and spatial relationships.'' So
%too would a method to churn butter or drill for oil. That is no
%accident: ``successions of bodily movement'' often serve basic functional
%purposes. Such movements do not become copyrightable as ``choreographic
%works'' when they are part and parcel of a process. Even if the Sequence
%could fit within some colloquial definitions of dance or choreography,
%it remains a process ineligible for copyright protection.

[The court rejected arguments that the sequence of poses was protectable as a
compilation or choreographic work. In particular:]
The idea/expression dichotomy, codified in Section 102(b), plays a
similar role in defining the scope of protection for a ``choreographic
work'' as it does for compilations.
In the context of choreographic works, that role is essential. Our
day-to-day lives consist of many routinized physical movements, from
brushing one's teeth to pushing a lawnmower to shaking a Polaroid
picture, that could be (and, in two of the preceding examples, have
been\readingfootnote{13}{\textit{See How To Do the `Lawn Mower' (Dance)},
WikiHow, \url{http://www.wikihow.com/Do-the-Lawn-Mower-(Dance)} (last visited
Sept. 16, 2015); \textit{Shake It Like a Polaroid Picture},
Urban Dictionary,
\url{http://www.urbandictionary.com/define.php?term=shake+it+like+a+Polaroid+picture
} (last visited Sept. 16, 2015).})
characterized as forms of dance. Without a proper understanding of the
idea/expression dichotomy, one might obtain monopoly rights over these
functional physical sequences by describing them in a tangible medium of
expression and labeling them choreographic works. The idea/expression
dichotomy thus ensures that expansive interpretations of the categories
enumerated as proper subjects of copyright will, ``[i]n no case,'' extend
copyright protection beyond its constitutional limits.


%
%\readinghead{IV. Conclusion}
%
%
%
%Although there is no cause to dispute the many health, fitness,
%spiritual, and aesthetic benefits of yoga, and Bikram Yoga in
%particular, they do not bring the Sequence into the realm of copyright
%protection. The Sequence falls squarely within Section 102(b)'s
%exclusions from copyright protection, no matter how it is labeled or how
%ably the label is argued. Therefore, the district court properly
%granted Evolation's motion for partial summary judgment.
%
%\textsc{Affirmed}.
%
